% GENERATED FILE. Edit canonical lessons, then run npm run generate:books.
% canonical-source-hash: fnv1a64:d5d9e63992dccb0c
% canonical-lessons: KA-C497-hogalu, KA-C497-abhimana, KA-C497-ahankara, KA-C497-vinaya, KA-C497-svabhimana, KA-R497-review-490-497

\chapter{Pride and Modesty}
\label{ch:ka-pride-and-modesty}
\hlchaptermodality{497}

\begin{chapteropening}
\textbf{By the end of this chapter:} \emph{I can say to praise, pride, arrogance, modesty and self-respect.}

Complete the last lesson of chapter 497: I can say to praise, pride, arrogance, modesty and self-respect.
\end{chapteropening}

\section[\texorpdfstring{\kn{ಹೊಗಳು} (hogaḷu)}{ಹೊಗಳು (hogaḷu)}]{\kn{ಹೊಗಳು} (hogaḷu) — to praise}

\label{lesson:KA-C497-hogalu}



\begin{quote}
Before the new one: say the Kannada for eagerness, then the Kannada for relieved.
\end{quote}

\subsection*{You'll want to know: \kn{ಹೊಗಳು}}
\textbf{\kn{ಹೊಗಳು}} — \emph{hogaḷu} — \textquotedblleft{}to praise\textquotedblright{}.

\begin{quote}
\kn{ಶಿಕ್ಷಕರು} \kn{ನನ್ನ} \kn{ಕಥೆಯನ್ನು} \kn{ಹೊಗಳಿದರು}.
\end{quote}

\begin{quote}
\emph{śikṣakaru nanna katheyannu hogaḷidaru.}
\end{quote}

\begin{quote}
The teacher praised my story.
\end{quote}

Past (I …): \kn{ಹೊಗಳಿದೆ}, \emph{hogaḷide}.

What is praised takes -\kn{ಅನ್ನು}. \kn{ಹೊಗಳಿಕೆ} is praise.

\begin{cousinweb}
\textbf{\kn{ಹೊಗಳು}} (\emph{hogaḷu}) is native Dravidian: DEDR 4235 lists an older Kannada form in p-, \textquotedblleft{}to praise; praise, renown\textquotedblright{}. Old p- became h- in the modern language.
\end{cousinweb}

\begin{grammarlens}[title={What you've built}]
One more word for describing what you lived and felt.
\end{grammarlens}

\subsection*{Your turn}
\begin{itemize}
\raggedright
  \item \emph{Say it:} \emph{hogaḷu}
  \item \emph{Say it:} \emph{hogaḷu}, once more
  \item \emph{Say it:} \emph{śikṣakaru nanna katheyannu hogaḷidaru.} — The teacher praised my story.
  \item \emph{Recall:} say the Kannada for eagerness, then the Kannada for relieved, then say \emph{hogaḷu} again
\end{itemize}

\begin{tcolorbox}[breakable,colback=teal!4,colframe=teal!35!black,title={Before you move on}]
What does \kn{ಹೊಗಳು} mean? (\textquotedblleft{}To praise\textquotedblright{}.) Say it once more.
\end{tcolorbox}
\section[\texorpdfstring{\kn{ಅಭಿಮಾನ} (abhimāna)}{ಅಭಿಮಾನ (abhimāna)}]{\kn{ಅಭಿಮಾನ} (abhimāna) — pride; fond regard}

\label{lesson:KA-C497-abhimana}



\begin{quote}
Before the new one: say the Kannada for relieved, then the Kannada for to praise.
\end{quote}

\subsection*{You'll want to know: \kn{ಅಭಿಮಾನ}}
\textbf{\kn{ಅಭಿಮಾನ}} — \emph{abhimāna} — \textquotedblleft{}pride; fond regard\textquotedblright{}.

\begin{quote}
\kn{ನಮ್ಮ} \kn{ಭಾಷೆಯ} \kn{ಬಗ್ಗೆ} \kn{ನಮಗೆ} \kn{ಅಭಿಮಾನ} \kn{ಇದೆ}.
\end{quote}

\begin{quote}
\emph{namma bhāṣeya bagge namage abhimāna ide.}
\end{quote}

\begin{quote}
We are proud of our language.
\end{quote}

Warm pride in what you love; \kn{ಅಭಿಮಾನಿ} is a fan.

\begin{cousinweb}
\textbf{\kn{ಅಭಿಮಾನ}} (\emph{abhimāna}) is borrowed from Sanskrit \emph{abhimāna}, \textquotedblleft{}pride\textquotedblright{} (CDIAL 524), from the root \emph{man}, \textquotedblleft{}to think\textquotedblright{}.
\end{cousinweb}

\begin{grammarlens}[title={What you've built}]
One more word for describing what you lived and felt.
\end{grammarlens}

\subsection*{Your turn}
\begin{itemize}
\raggedright
  \item \emph{Say it:} \emph{abhimāna}
  \item \emph{Say it:} \emph{abhimāna}, once more
  \item \emph{Say it:} \emph{namma bhāṣeya bagge namage abhimāna ide.} — We are proud of our language.
  \item \emph{Recall:} say the Kannada for relieved, then the Kannada for to praise, then say \emph{abhimāna} again
\end{itemize}

\begin{tcolorbox}[breakable,colback=teal!4,colframe=teal!35!black,title={Before you move on}]
What does \kn{ಅಭಿಮಾನ} mean? (\textquotedblleft{}Pride; fond regard\textquotedblright{}.) Say it once more.
\end{tcolorbox}
\section[\texorpdfstring{\kn{ಅಹಂಕಾರ} (ahaṅkāra)}{ಅಹಂಕಾರ (ahaṅkāra)}]{\kn{ಅಹಂಕಾರ} (ahaṅkāra) — arrogance; conceit}

\label{lesson:KA-C497-ahankara}



\begin{quote}
Before the new one: say the Kannada for to praise, then the Kannada for pride.
\end{quote}

\subsection*{You'll want to know: \kn{ಅಹಂಕಾರ}}
\textbf{\kn{ಅಹಂಕಾರ}} — \emph{ahaṅkāra} — \textquotedblleft{}arrogance; conceit\textquotedblright{}.

\begin{quote}
\kn{ಹಣ} \kn{ಬಂದ} \kn{ಮೇಲೆ} \kn{ಅವನಿಗೆ} \kn{ಅಹಂಕಾರ} \kn{ಬಂದಿತು}.
\end{quote}

\begin{quote}
\emph{haṇa banda mēle avanige ahaṅkāra banditu.}
\end{quote}

\begin{quote}
Once the money came, he grew arrogant.
\end{quote}

From \kn{ಅಹಂ}, I, and \kn{ಕಾರ}, making: making everything about \textquotedblleft{}I\textquotedblright{}.

\begin{cousinweb}
\textbf{\kn{ಅಹಂಕಾರ}} (\emph{ahaṅkāra}) is borrowed from Sanskrit \emph{ahaṃkāra}, \textquotedblleft{}the making of self, egotism, pride\textquotedblright{} (Monier-Williams).
\end{cousinweb}

\begin{grammarlens}[title={What you've built}]
One more word for describing what you lived and felt.
\end{grammarlens}

\subsection*{Your turn}
\begin{itemize}
\raggedright
  \item \emph{Say it:} \emph{ahaṅkāra}
  \item \emph{Say it:} \emph{ahaṅkāra}, once more
  \item \emph{Say it:} \emph{haṇa banda mēle avanige ahaṅkāra banditu.} — Once the money came, he grew arrogant.
  \item \emph{Recall:} say the Kannada for to praise, then the Kannada for pride, then say \emph{ahaṅkāra} again
\end{itemize}

\begin{tcolorbox}[breakable,colback=teal!4,colframe=teal!35!black,title={Before you move on}]
What does \kn{ಅಹಂಕಾರ} mean? (\textquotedblleft{}Arrogance; conceit\textquotedblright{}.) Say it once more.
\end{tcolorbox}
\section[\texorpdfstring{\kn{ವಿನಯ} (vinaya)}{ವಿನಯ (vinaya)}]{\kn{ವಿನಯ} (vinaya) — modesty; humility}

\label{lesson:KA-C497-vinaya}



\begin{quote}
Before the new one: say the Kannada for pride, then the Kannada for arrogance.
\end{quote}

\subsection*{You'll want to know: \kn{ವಿನಯ}}
\textbf{\kn{ವಿನಯ}} — \emph{vinaya} — \textquotedblleft{}modesty; humility\textquotedblright{}.

\begin{quote}
\kn{ದೊಡ್ಡ} \kn{ಬಹುಮಾನ} \kn{ಗೆದ್ದರೂ} \kn{ಅವಳಲ್ಲಿ} \kn{ವಿನಯ} \kn{ಇದೆ}.
\end{quote}

\begin{quote}
\emph{doḍḍa bahumāna geddarū avaḷalli vinaya ide.}
\end{quote}

\begin{quote}
Though she won a big prize, she is modest.
\end{quote}

\kn{ವಿನಯದಿಂದ}, politely, humbly.

\begin{cousinweb}
\textbf{\kn{ವಿನಯ}} (\emph{vinaya}) is borrowed from Sanskrit \emph{vinaya}, \textquotedblleft{}leading away; guidance, good breeding\textquotedblright{} (CDIAL 11767).
\end{cousinweb}

\begin{grammarlens}[title={What you've built}]
One more word for describing what you lived and felt.
\end{grammarlens}

\subsection*{Your turn}
\begin{itemize}
\raggedright
  \item \emph{Say it:} \emph{vinaya}
  \item \emph{Say it:} \emph{vinaya}, once more
  \item \emph{Say it:} \emph{doḍḍa bahumāna geddarū avaḷalli vinaya ide.} — Though she won a big prize, she is modest.
  \item \emph{Recall:} say the Kannada for pride, then the Kannada for arrogance, then say \emph{vinaya} again
\end{itemize}

\begin{tcolorbox}[breakable,colback=teal!4,colframe=teal!35!black,title={Before you move on}]
What does \kn{ವಿನಯ} mean? (\textquotedblleft{}Modesty; humility\textquotedblright{}.) Say it once more.
\end{tcolorbox}
\section[\texorpdfstring{\kn{ಸ್ವಾಭಿಮಾನ} (svābhimāna)}{ಸ್ವಾಭಿಮಾನ (svābhimāna)}]{\kn{ಸ್ವಾಭಿಮಾನ} (svābhimāna) — self-respect}

\label{lesson:KA-C497-svabhimana}



\begin{quote}
Before the new one: say the Kannada for arrogance, then the Kannada for modesty.
\end{quote}

\subsection*{You'll want to know: \kn{ಸ್ವಾಭಿಮಾನ}}
\textbf{\kn{ಸ್ವಾಭಿಮಾನ}} — \emph{svābhimāna} — \textquotedblleft{}self-respect\textquotedblright{}.

\begin{quote}
\kn{ಅವನು} \kn{ಸ್ವಾಭಿಮಾನದಿಂದ} \kn{ಆ} \kn{ಹಣವನ್ನು} \kn{ನಿರಾಕರಿಸಿದನು}.
\end{quote}

\begin{quote}
\emph{avanu svābhimānadinda ā haṇavannu nirākarisidanu.}
\end{quote}

\begin{quote}
Out of self-respect he refused that money.
\end{quote}

\kn{ಸ್ವ}, self + \kn{ಅಭಿಮಾನ}, pride: the word from earlier in this chapter, turned on yourself.

\begin{cousinweb}
\textbf{\kn{ಸ್ವಾಭಿಮಾನ}} (\emph{svābhimāna}) is a compound built from Sanskrit: \emph{sva}, \textquotedblleft{}own, self\textquotedblright{}, and \emph{abhimāna}, \textquotedblleft{}pride\textquotedblright{} (CDIAL 524).
\end{cousinweb}

\begin{grammarlens}[title={What you've built}]
One more word for describing what you lived and felt.
\end{grammarlens}

\subsection*{Your turn}
\begin{itemize}
\raggedright
  \item \emph{Say it:} \emph{svābhimāna}
  \item \emph{Say it:} \emph{svābhimāna}, once more
  \item \emph{Say it:} \emph{avanu svābhimānadinda ā haṇavannu nirākarisidanu.} — Out of self-respect he refused that money.
  \item \emph{Recall:} say the Kannada for arrogance, then the Kannada for modesty, then say \emph{svābhimāna} again
\end{itemize}

\begin{tcolorbox}[breakable,colback=teal!4,colframe=teal!35!black,title={Before you move on}]
What does \kn{ಸ್ವಾಭಿಮಾನ} mean? (\textquotedblleft{}Self-respect\textquotedblright{}.) Say it once more.
\end{tcolorbox}
\section[(dialogue)]{Looking back: fear, grief, joy and pride}

\label{lesson:KA-R497-review-490-497}



\begin{quote}
Say the Kannada for to be afraid and to feel shy.
\end{quote}

\subsection*{Your turn}
Say these aloud:

\begin{itemize}
\raggedright
  \item to be afraid, to feel shy, anxiety, embarrassment, unease
  \item to pine, to feel pity, anguish, regret, remorse
  \item cheerfulness, to enjoy, zest, delight, satisfaction
  \item to get annoyed, to provoke, to ridicule, frustration, agitation
  \item to shrink back, to fumble, amazement, to lose heart, bewilderment
  \item to cuddle, to kiss, tender affection, motherly love, care
  \item to long for, to rue a loss, hope, eagerness, relieved
  \item to praise, pride, arrogance, modesty, self-respect
\end{itemize}

\begin{tcolorbox}[breakable,colback=teal!4,colframe=teal!35!black,title={Before you move on}]
To be afraid? (\textbf{\kn{ಹೆದರು}}, \emph{hedaru}.) To feel shy? (\textbf{\kn{ನಾಚು}}, \emph{nācu}.) Anxiety? (\textbf{\kn{ಆತಂಕ}}, \emph{ātaṅka}.) Embarrassment? (\textbf{\kn{ಮುಜುಗರ}}, \emph{mujugara}.) Unease? (\textbf{\kn{ಕಳವಳ}}, \emph{kaḷavaḷa}.) To pine? (\textbf{\kn{ಕೊರಗು}}, \emph{koragu}.) To feel pity? (\textbf{\kn{ಮರುಗು}}, \emph{marugu}.) Anguish? (\textbf{\kn{ಸಂಕಟ}}, \emph{saṅkaṭa}.) Regret? (\textbf{\kn{ವಿಷಾದ}}, \emph{viṣāda}.) Remorse? (\textbf{\kn{ಪಶ್ಚಾತ್ತಾಪ}}, \emph{paścāttāpa}.) Cheerfulness? (\textbf{\kn{ಉಲ್ಲಾಸ}}, \emph{ullāsa}.) To enjoy? (\textbf{\kn{ಆನಂದಿಸು}}, \emph{ānandisu}.) Zest? (\textbf{\kn{ಹುಮ್ಮಸ್ಸು}}, \emph{hummassu}.) Delight? (\textbf{\kn{ಹರ್ಷ}}, \emph{harṣa}.) Satisfaction? (\textbf{\kn{ತೃಪ್ತಿ}}, \emph{tṛpti}.) To get annoyed? (\textbf{\kn{ಸಿಟ್ಟಾಗು}}, \emph{siṭṭāgu}.) To provoke? (\textbf{\kn{ಕೆಣಕು}}, \emph{keṇaku}.) To ridicule? (\textbf{\kn{ಹೀಯಾಳಿಸು}}, \emph{hīyāḷisu}.) Frustration? (\textbf{\kn{ಹತಾಶೆ}}, \emph{hatāśe}.) Agitation? (\textbf{\kn{ಉದ್ವೇಗ}}, \emph{udvēga}.) To shrink back? (\textbf{\kn{ಮುದುಡು}}, \emph{muduḍu}.) To fumble? (\textbf{\kn{ತಡವರಿಸು}}, \emph{taḍavarisu}.) Amazement? (\textbf{\kn{ಬೆರಗು}}, \emph{beragu}.) To lose heart? (\textbf{\kn{ಕಂಗೆಡು}}, \emph{kaṅgeḍu}.) Bewilderment? (\textbf{\kn{ದಿಗ್ಭ್ರಮೆ}}, \emph{digbhrame}.) To cuddle? (\textbf{\kn{ಮುದ್ದಿಸು}}, \emph{muddisu}.) To kiss? (\textbf{\kn{ಮುತ್ತಿಡು}}, \emph{muttiḍu}.) Tender affection? (\textbf{\kn{ವಾತ್ಸಲ್ಯ}}, \emph{vātsalya}.) Motherly love? (\textbf{\kn{ಮಮತೆ}}, \emph{mamate}.) Care? (\textbf{\kn{ಕಾಳಜಿ}}, \emph{kāḷaji}.) To long for? (\textbf{\kn{ಹಂಬಲಿಸು}}, \emph{hambalisu}.) To rue a loss? (\textbf{\kn{ಹಳಹಳಿಸು}}, \emph{haḷahaḷisu}.) Hope? (\textbf{\kn{ಭರವಸೆ}}, \emph{bharavase}.) Eagerness? (\textbf{\kn{ಕಾತರ}}, \emph{kātara}.) Relieved? (\textbf{\kn{ನಿರಾಳ}}, \emph{nirāḷa}.) To praise? (\textbf{\kn{ಹೊಗಳು}}, \emph{hogaḷu}.) Pride? (\textbf{\kn{ಅಭಿಮಾನ}}, \emph{abhimāna}.) Arrogance? (\textbf{\kn{ಅಹಂಕಾರ}}, \emph{ahaṅkāra}.) Modesty? (\textbf{\kn{ವಿನಯ}}, \emph{vinaya}.) Self-respect? (\textbf{\kn{ಸ್ವಾಭಿಮಾನ}}, \emph{svābhimāna}.)
\end{tcolorbox}
